\begin{table}[!htbp]
\centering
\caption{Representative forecasting studies at intra-hour and short
hourly horizons. The studies use different targets and test protocols, so
their results are not directly ranked.}
\label{tab:literature_review}
\label{tab:literature_review_short}
\scriptsize
\setlength{\tabcolsep}{2.5pt}
\renewcommand{\arraystretch}{1.12}
\begin{tabular}{@{}p{0.11\linewidth}p{0.18\linewidth}p{0.22\linewidth}p{0.18\linewidth}p{0.25\linewidth}@{}}
\toprule
\textbf{Study} & \textbf{Target and scope} & \textbf{Inputs and processing} &
\textbf{Predictor} & \textbf{Horizon and evaluation design} \\
\midrule
Chow et al.\ \cite{chow2011} &
GHI and cloud conditions at the UC San Diego test bed &
30-s total-sky images, clear-sky library, cloud map, shadows, and motion vectors &
Image-based physical nowcast &
Cloud nowcasts up to 5~min; six pyranometers on four partly cloudy days; qualitative GHI assessment \\

Guariso et al.\ \cite{guariso2020} &
Hourly GHI at an Italian station, with transfer to other sites &
Autoregressive irradiance windows; no decomposition &
Feed-forward and LSTM networks &
1--12~h; recursive, horizon-specific, and multi-output strategies; persistence and clear-sky references \\

Huang et al.\ \cite{huang2021} &
Hourly irradiance at Denver, Clark, and Folsom &
Irradiance history and meteorology, including next-hour temperature, humidity, and wind-speed forecasts; wavelet-packet decomposition &
WPD--CNN--LSTM--MLP hybrid &
1~h ahead; comparison with individual networks and climatology--persistence \\
\bottomrule
\end{tabular}
\end{table}

\begin{table}[!htbp]
\centering
\caption{Representative day-ahead and direct multi-step solar-forecasting
studies.}
\label{tab:literature_review_multistep}
\scriptsize
\setlength{\tabcolsep}{2.5pt}
\renewcommand{\arraystretch}{1.12}
\begin{tabular}{@{}p{0.11\linewidth}p{0.18\linewidth}p{0.22\linewidth}p{0.18\linewidth}p{0.25\linewidth}@{}}
\toprule
\textbf{Study} & \textbf{Target and scope} & \textbf{Inputs and processing} &
\textbf{Predictor} & \textbf{Horizon and evaluation design} \\
\midrule
Cheng et al.\ \cite{cheng2021} &
Daytime hourly irradiance in Taiwan under tropical and subtropical marine conditions &
Previous days of irradiance and selected temperature, humidity, and rainfall observations or forecasts &
One-dimensional CNN, LSTM, skip concatenation, and dense output &
Day-ahead through 7~days ahead; limited-data deployment experiment \\

Voyant et al.\ \cite{voyant2022} &
Solar-radiation series from sites in multiple climates &
Little or no training history; statistical structure of the target &
Five na\"ive references (including a second-order autoregressive model) and their forecast combination &
Direct multi-step benchmark evaluated with several complementary metrics \\

Zhao et al.\ \cite{zhao2024timesnet} &
Hourly solar radiation &
ICEEMDAN-decomposed solar-radiation sequences combined with meteorological data &
ICEEMDAN--TimesNet &
Week-ahead multi-step prediction; comparison with alternative multi-step models \\
\bottomrule
\end{tabular}
\end{table}

\begin{table}[!htbp]
\centering
\caption{Representative monthly, multi-climate, and very-long-horizon
forecasting studies.}
\label{tab:literature_review_long}
\scriptsize
\setlength{\tabcolsep}{2.5pt}
\renewcommand{\arraystretch}{1.12}
\begin{tabular}{@{}p{0.11\linewidth}p{0.18\linewidth}p{0.22\linewidth}p{0.18\linewidth}p{0.25\linewidth}@{}}
\toprule
\textbf{Study} & \textbf{Target and scope} & \textbf{Inputs and processing} &
\textbf{Predictor} & \textbf{Horizon and evaluation design} \\
\midrule
Ozoegwu \cite{ozoegwu2019} &
Monthly-mean daily global solar radiation at selected locations &
Radiation history; lagged month numbers in NARX and current month number in the hybrid &
NAR, NARX, and hybrid ANN &
Monthly-scale forecasting; comparison among the three neural formulations \\

Zhu et al.\ \cite{zhu2024} &
Solar irradiance in a photovoltaic-greenhouse setting &
Multiple meteorological-factor combinations and training-set lengths &
TimesNet, Informer, Transformer, and Pyraformer &
Nine horizons from 10~min to 24~h; seasonal and full-climate analyses \\

Cabral et al.\ \cite{cabral2026} &
GHI at 40 locations spanning multiple climate classes &
GHI-only enhanced pattern-recognition pipeline &
GPT-Neo and contemporary neural competitors &
Very-long-term multi-year prediction; location and aggregate errors plus computational assessment \\
\bottomrule
\end{tabular}
\par\medskip
\begin{minipage}{0.96\linewidth}
\scriptsize
\textit{Abbreviations in Tables~\ref{tab:literature_review_short}--%
\ref{tab:literature_review_long}:} ANN, artificial neural network; CNN,
convolutional neural network; GHI, global horizontal irradiance; ICEEMDAN,
improved complete ensemble empirical mode decomposition with adaptive noise;
LSTM, long short-term memory; MLP, multilayer perceptron; NAR, nonlinear
autoregressive; NARX, nonlinear autoregressive with exogenous inputs; WPD,
wavelet packet decomposition.
\end{minipage}
\end{table}
